\section{Reflection}

This was the most hands-on assignment I have had in the course so far. Getting something
to actually run in a real cluster, with real constraints, is very different from following
a tutorial where everything is set up for you.

The local containerization phase was where I made the most mistakes, but also where I
learned the most concretely. The architecture mismatch between ARM64 and AMD64 is
something I would not have run into if I had been on a standard x86 machine, but having
to debug an \texttt{Illegal Instruction} crash taught me more about how Go cross-compilation
works than any amount of reading would have. The UBI permission issue was similar. Knowing
in theory that containers should run as non-root is different from actually seeing the
permission error, tracing it to \texttt{/app} being root-owned, and finding
\texttt{/opt/app-root/src} as the right answer. Those kinds of problems leave a stronger
impression than theory does.

The ArgoCD section was frustrating but valuable. I came in expecting to follow a
relatively straightforward installation guide and instead spent several sessions debugging
RBAC errors and discovering the edges of the Sandbox environment. By the end of it I had
a much clearer picture of what ArgoCD actually needs to operate than I would have gotten
from a successful install. A tool that works perfectly does not teach you what it depends
on. A tool that fails in twelve different ways does.

I also came away with a better understanding of why the two-repository pattern exists.
It felt like extra ceremony at first, maintaining a separate config repo just to store
YAML files. But once I had the GitHub Actions workflow updating the image tag and could
see exactly how the code side and the infrastructure side are kept separate, it started
to make sense. In a team environment where different people have access to different
systems, that separation is not just organization, it is a security boundary.

If I were to do this again I would spend more time reading the base image documentation
before writing the first Containerfile. Most of the early errors came from assumptions
that turned out to be wrong, and the documentation for the Red Hat UBI images is quite
clear about the non-root conventions once you know to look for it.
